\scanpage{76}

\deleted{(\unsure{enfin} la surjectivité est \unsure{acquise}) que $T_\alpha,T_\beta$ sont dans une $=$ comp. connexe de $X$, il existe une ``cha\unsure{î|i}ne'' d'\uline{\unsure{indices}}}

\[
\deleted{i_0=i,\ i_1,\ \mathord{-}\ \mathord{-},\ i_n=j}
\]

\deleted{telle que pour deux indices consécutifs $k=i_{\unsure{r|\nu}}$, $k'=i_{\unsure{r|\nu}+1}$, on ait $T_k\subset T_{k'}$ ou $T_{k'}\subset T_k$.}

\margin{Supposons de plus que les intérieurs des $T_\alpha$ recouvrent $X$ \unsure{ou} que les $T_\alpha$ soient fermés (\deleted{\unsure{et}} $A$ fini) Alors \unsure{on a} $\Pi_0(A)\xrightarrow{\sim}\Pi_0(X)$. Mieux ---}

Supposons de plus \inserted{que les intérieurs des $T_i$ recouvrent $X$,} \deleted{\unsure{que les $T_i$ s'}} \inserted{\uline{Ceci posé,}} je dis que la connaissance des (pour $u$ fixé) des $g_{y,\unsure{x}}^\alpha(u)\in\pi_\alpha$ \inserted{($\alpha\in A,y,x\in T_\alpha$)} permet de retrouver tous les invariants $g_l(u)$ (ou $f_l(u)$). En effet, il suffit de voir que les flèches $l_{y,x}^\alpha=p_\alpha(l_{y,x})$ ($\alpha\in I$, $x,y\in T_\alpha$, $p_\alpha:T_\alpha\hookrightarrow X$ l'inclusion) \deleted{\unsure{eng}} forment un syst. de générateurs du groupoïde $\Pi_1(X_K)\deleted{\unsure{x_0})}$ --- \inserted{tous} des \deleted{obj} sommets sont obtenus en tous cas, et on trouve une application bijective sur les $\Pi_0$ (du groupoïde engendré par les gén. $l_{y,x}^\alpha$ \inserted{($\alpha\in I,x,y\in T_{\alpha K}$)} \unsure{\inserted{et les relations $l_{y,x}^{\unsure{\beta}}=\varphi_{\beta\alpha}(l_{y,x}^{\alpha})$}} $\unsure{\Big\}}$ et de $\Pi_1(X_K)$), il faut voir qu'on a une application surjective sur les $\Pi_1$.

Quand les intérieurs des $T_{\unsure{\alpha}}$ recouvrent $X$, il suffit de noter que la connaissance des $g_{\unsure{y},x}^{\alpha}(u)$ nous donne les $g_{y,x}(u)$ (sans exposant $\alpha$ !) pour $x,y$ points ``voisins'' sur $X$,

\margin{75}
